\documentclass[11pt,a4paper]{article}
\usepackage[margin=1in]{geometry}
\usepackage{amsmath,amssymb,mathtools}
\usepackage{enumitem}
\usepackage{hyperref}
\usepackage[T1]{fontenc}
\usepackage{lmodern}
\setlength{\parindent}{0pt}
\setlength{\parskip}{0.65em}

\newcommand{\R}{\mathbb{R}}
\newcommand{\C}{\mathbb{C}}
\newcommand{\K}{\mathbb{K}}
\newcommand{\Span}{\operatorname{span}}
\newcommand{\rank}{\operatorname{rank}}

\title{Linear Algebra -- Problem Set 10}
\author{Taras Halay}
\date{}

\begin{document}
\maketitle

\section*{Exercise 3.17}

\subsection*{Statement summary}
Let $E=\R_3[X]$ be the vector space of real polynomials of degree at most $3$. Consider the symmetric bilinear form
\[
    \omega(p,q)=\int_{-1}^{1}p(x)q(x)\,dx.
\]
Using the standard basis $\mathcal E=(1,x,x^2,x^3)$, compute the matrix of $\omega$. Then use the induced inner-product geometry to approximate $R(x)=x^3$ by a polynomial from $F=\R_2[X]$ in the least-squares sense.

\subsection*{Solution}
\subsubsection*{(i) Matrix of the bilinear form}
Let $\mathcal E=\{1,x,x^2,x^3\}$ be the standard basis of $E$. Then
\[
[\omega]^{\mathcal E}_{\mathcal E}=\begin{pmatrix}
\omega(1,1)&\omega(1,x)&\omega(1,x^2)&\omega(1,x^3)\\
\omega(x,1)&\omega(x,x)&\omega(x,x^2)&\omega(x,x^3)\\
\omega(x^2,1)&\omega(x^2,x)&\omega(x^2,x^2)&\omega(x^2,x^3)\\
\omega(x^3,1)&\omega(x^3,x)&\omega(x^3,x^2)&\omega(x^3,x^3)
\end{pmatrix}.
\]
For $1\le i,j\le 4$,
\[
([\omega]^{\mathcal E}_{\mathcal E})_{ij}
=\omega(x^{i-1},x^{j-1})
=\int_{-1}^{1}x^{i+j-2}\,dx
=\left.\frac{x^{i+j-1}}{i+j-1}\right|_{-1}^{1}.
\]
The integral is zero when $i+j-2$ is odd, equivalently when $i+j$ is odd. If $i+j$ is even, then
\[
\int_{-1}^{1}x^{i+j-2}\,dx=\frac{2}{i+j-1}.
\]
Therefore
\[
[\omega]^{\mathcal E}_{\mathcal E}=\begin{pmatrix}
2&0&\frac23&0\\
0&\frac23&0&\frac25\\
\frac23&0&\frac25&0\\
0&\frac25&0&\frac27
\end{pmatrix}.
\]

\subsubsection*{(ii) Least-squares functional}
Let $R(x)=x^3$. For $Q\in F=\R_2[X]$,
\[
J(Q)=\omega(x^3-Q(x),x^3-Q(x))=\|x^3-Q(x)\|_\omega^2.
\]
Thus minimizing $J$ is exactly the problem of finding the closest point to $x^3$ inside the subspace $F$ with respect to the norm induced by $\omega$.

\subsubsection*{(iii) Orthonormal basis of $F$}
We apply the Gram-Schmidt process to the ordered basis $(1,x,x^2)$ of $F$.
First,
\[
u_1=\frac{1}{\|1\|_\omega}=\frac{1}{\sqrt2}.
\]
Next,
\[
u_2=\frac{x-\omega(x,u_1)u_1}{\|x-\omega(x,u_1)u_1\|_\omega}.
\]
Since $\omega(x,u_1)=0$, we get
\[
u_2=\frac{x}{\|x\|_\omega}=\frac{x}{\sqrt{2/3}}=\frac{\sqrt3}{\sqrt2}x.
\]
For the third vector,
\[
u_3=\frac{x^2-\sum_{i=1}^{2}\omega(x^2,u_i)u_i}{\left\|x^2-\sum_{i=1}^{2}\omega(x^2,u_i)u_i\right\|_\omega}.
\]
We have
\[
\omega(x^2,u_1)u_1=\left(\int_{-1}^{1}x^2\frac{1}{\sqrt2}\,dx\right)\frac{1}{\sqrt2}=\frac13,
\qquad
\omega(x^2,u_2)u_2=0.
\]
Hence
\[
u_3=\frac{x^2-\frac13}{\left\|x^2-\frac13\right\|_\omega}.
\]
Moreover,
\begin{align*}
\left\|x^2-\frac13\right\|_\omega^2
&=\omega\left(x^2-\frac13,x^2-\frac13\right)\\
&=\int_{-1}^{1}\left(x^2-\frac13\right)^2dx\\
&=\int_{-1}^{1}\left(x^4-\frac23x^2+\frac19\right)dx\\
&=\frac25-\frac49+\frac29=\frac{8}{45}.
\end{align*}
Thus an orthonormal basis of $F$ is
\[
\mathcal U=\left\{\frac1{\sqrt2},\frac{\sqrt3}{\sqrt2}x,
\frac{x^2-\frac13}{\sqrt{8/45}}\right\}.
\]

\subsubsection*{(iv) Projection of $x^3$ onto $F$}
Since $\mathcal U$ is orthonormal,
\[
P_F(R)=\sum_{i=1}^{3}\langle R,u_i\rangle u_i.
\]
We note that
\[
\langle x^3,u_1\rangle=0,\qquad \langle x^3,u_3\rangle=0,
\]
because the corresponding integrands are odd. Therefore
\[
P_F(R)=\omega\left(x^3,\frac{\sqrt3}{\sqrt2}x\right)\frac{\sqrt3}{\sqrt2}x.
\]
So
\[
P_F(R)=\frac{\sqrt3}{\sqrt2}\int_{-1}^{1}x^4\,dx\cdot\frac{\sqrt3}{\sqrt2}x
=\frac32\cdot\frac25x=\frac35x.
\]
The minimal value is
\begin{align*}
\omega\left(x^3-\frac35x,x^3-\frac35x\right)
&=\int_{-1}^{1}\left(x^3-\frac35x\right)^2dx\\
&=\int_{-1}^{1}\left(x^6-\frac65x^4+\frac9{25}x^2\right)dx\\
&=\frac27-\frac65\cdot\frac25+\frac9{25}\cdot\frac23
=\frac{8}{175}.
\end{align*}

The conceptual reason is least-squares geometry: if $Q^*$ is the best approximation of $x^3$ in $F$, then
\[
x^3-Q^*\perp F.
\]
Equivalently,
\[
\langle x^3-Q^*,S\rangle=0,\qquad \forall S\in\R_2[X].
\]
Therefore the closest polynomial is the orthogonal projection of $x^3$ onto $F$, namely
\[
Q^*(x)=\frac35x.
\]

\newpage
\section*{Exercise 4.7}

\subsection*{Statement summary}
Let $x=(x_0,x_1,\dots,x_m)\in\K^{m+1}$ and define the Vandermonde matrix
\[
V(x)=
\begin{pmatrix}
1&x_0&x_0^2&\cdots&x_0^n\\
1&x_1&x_1^2&\cdots&x_1^n\\
1&x_2&x_2^2&\cdots&x_2^n\\
\vdots&\vdots&\vdots&\ddots&\vdots\\
1&x_m&x_m^2&\cdots&x_m^n
\end{pmatrix}.
\]
In the square case $m=n$, prove the vanishing criterion for $\det V(x)$, show that each difference $x_i-x_j$ divides the determinant, and derive the closed formula for the Vandermonde determinant.

\subsection*{Solution}
In the following parts, $n=m$.

\subsubsection*{(i) Vanishing criterion}
We prove that
\[
\det V(x)=0 \iff x_i=x_j\quad\text{for some }i\ne j.
\]
Assume first that $x_i=x_j$ for some $i\ne j$. Then row $R_i$ equals row $R_j$:
\[
R_i=(1,x_i,x_i^2,\dots,x_i^n)=(1,x_j,x_j^2,\dots,x_j^n)=R_j.
\]
Since the determinant is alternating in the rows, $\det V(x)=0$.

Conversely, assume $\det V(x)=0$. Then $\rank V(x)<n+1$, so the columns are linearly dependent. Hence there exists
\[
(a_0,\dots,a_n)^T\ne 0
\]
such that
\[
V(x)\begin{pmatrix}a_0\\ \vdots\\ a_n\end{pmatrix}=0.
\]
Thus, for $r=0,\dots,n$,
\[
a_0+a_1x_r+\cdots+a_nx_r^n=0.
\]
Define
\[
p(t)=a_0+a_1t+\cdots+a_nt^n.
\]
This polynomial is not the zero polynomial, and it has roots $x_0,\dots,x_n$. Since $\deg p\le n$, it cannot have $n+1$ distinct roots. Therefore
\[
x_i=x_j
\]
for some $i\ne j$.

\subsubsection*{(ii) Divisibility by pairwise differences}
Fix $0\le i<j\le n$. If $x_i=x_j$, then by part (i),
\[
\det V(x)=0.
\]
Consider $\det V(x)$ as a polynomial in the variable $x_i$, with the other variables fixed. Since it vanishes when $x_i=x_j$, the factor theorem gives
\[
x_i-x_j \mid \det V(x).
\]
Equivalently, up to sign,
\[
x_j-x_i\mid \det V(x).
\]
Thus every pairwise difference divides the determinant.

\subsubsection*{(iii) Vandermonde determinant formula}
By the Leibniz formula,
\[
\det V(x)=\sum_{\sigma\in S_{n+1}}\operatorname{sgn}(\sigma)\prod_{r=0}^{n}x_r^{\sigma(r)}.
\]
Each term has the form
\[
x_0^{\sigma(0)}x_1^{\sigma(1)}\cdots x_n^{\sigma(n)},
\]
and its total degree is
\[
\sum_{r=0}^{n}\sigma(r)=0+1+\cdots+n=\frac{n(n+1)}2.
\]
Therefore
\[
\deg(\det V(x))\le \frac{n(n+1)}2.
\]
Now inspect the degree of
\[
\prod_{0\le i<j\le n}(x_j-x_i).
\]
There are $\binom{n+1}{2}=\frac{n(n+1)}2$ factors, each of degree $1$. Hence
\[
\deg\left(\prod_{0\le i<j\le n}(x_j-x_i)\right)=\frac{n(n+1)}2.
\]
From part (ii), the product divides $\det V(x)$. Since both polynomials have the same degree, there is a constant $C\in\K$ such that
\[
\det V(x)=C\prod_{0\le i<j\le n}(x_j-x_i).
\]
It remains to identify $C$. Compare the coefficient of the monomial
\[
x_0^0x_1^1\cdots x_n^n.
\]
In the Leibniz formula, this term comes from the identity permutation, whose sign is $1$. Hence its coefficient in $\det V(x)$ is $1$.

In the product
\[
\prod_{0\le i<j\le n}(x_j-x_i),
\]
we obtain the same monomial by choosing $x_j$ from every factor. Each chosen term has coefficient $1$, so the coefficient of $x_0^0x_1^1\cdots x_n^n$ is also $1$. Thus $C=1$ and
\[
\boxed{\det V(x)=\prod_{0\le i<j\le n}(x_j-x_i)}.
\]

\newpage
\section*{Exercise 5.1}

\subsection*{Statement summary}
For the two matrices
\[
A=\begin{pmatrix}
3&0&-1\\
2&4&2\\
-1&0&3
\end{pmatrix},
\qquad
B=\begin{pmatrix}
0&0&-8\\
0&2&0\\
4&0&12
\end{pmatrix},
\]
find the characteristic polynomials, eigenvalues, eigenspaces, and decide whether the matrices are diagonalizable. When possible, use diagonalization to compute $A^n$ and $B^n$.

\subsection*{Solution}
\subsection*{Matrix $A$}

\subsubsection*{(i) Characteristic polynomial}
We compute
\[
A-\lambda I=\begin{pmatrix}
3-\lambda&0&-1\\
2&4-\lambda&2\\
-1&0&3-\lambda
\end{pmatrix}.
\]
Then
\begin{align*}
\det(A-\lambda I)
&=(4-\lambda)\det\begin{pmatrix}3-\lambda&-1\\-1&3-\lambda\end{pmatrix}\\
&=(4-\lambda)((3-\lambda)^2-1)\\
&=(4-\lambda)^2(2-\lambda).
\end{align*}
Hence the eigenvalues are
\[
\lambda_1=4,\qquad \lambda_2=2,
\]
with algebraic multiplicities
\[
m_4=2,\qquad m_2=1.
\]

\subsubsection*{(ii) Eigenspaces and diagonalizability}
For $\lambda=4$,
\[
V_4^A=\ker(A-4I)=\ker\begin{pmatrix}
-1&0&-1\\
2&0&2\\
-1&0&-1
\end{pmatrix}.
\]
The equations give $x+z=0$. Therefore
\[
V_4^A=\Span\left\{\begin{pmatrix}1\\0\\-1\end{pmatrix},
\begin{pmatrix}0\\1\\0\end{pmatrix}\right\}.
\]
For $\lambda=2$,
\[
V_2^A=\ker(A-2I)=\ker\begin{pmatrix}
1&0&-1\\
2&2&2\\
-1&0&1
\end{pmatrix}.
\]
We obtain
\[
x-z=0,\qquad x+y+z=0.
\]
Thus $x=z$ and $y=-2x$, so
\[
V_2^A=\Span\left\{\begin{pmatrix}1\\-2\\1\end{pmatrix}\right\}.
\]
Here
\[
m_2=d_2=1,
\qquad
m_4=d_4=2.
\]
Hence $A$ is diagonalizable.

\subsubsection*{(iii) Computing $A^n$}
Let
\[
u_1=\begin{pmatrix}1\\-2\\1\end{pmatrix},
\qquad
u_2=\begin{pmatrix}1\\0\\-1\end{pmatrix},
\qquad
u_3=\begin{pmatrix}0\\1\\0\end{pmatrix}.
\]
Then
\[
P=(u_1,u_2,u_3)=\begin{pmatrix}
1&1&0\\
-2&0&1\\
1&-1&0
\end{pmatrix},
\qquad
D=\begin{pmatrix}
2&0&0\\
0&4&0\\
0&0&4
\end{pmatrix}.
\]
Also
\[
P^{-1}=\begin{pmatrix}
\frac12&0&\frac12\\
\frac12&0&-\frac12\\
1&1&1
\end{pmatrix}.
\]
Since $A=PDP^{-1}$, we have
\[
A^n=PD^nP^{-1}.
\]
Therefore
\[
A^n=
\begin{pmatrix}
\frac{2^n+4^n}{2}&0&\frac{2^n-4^n}{2}\\
-2^n+4^n&4^n&-2^n+4^n\\
\frac{2^n-4^n}{2}&0&\frac{2^n+4^n}{2}
\end{pmatrix}.
\]

\subsection*{Matrix $B$}
\subsubsection*{(i) Characteristic polynomial}
We compute
\[
B-\lambda I=\begin{pmatrix}
-\lambda&0&-8\\
0&2-\lambda&0\\
4&0&12-\lambda
\end{pmatrix}.
\]
Thus
\begin{align*}
\det(B-\lambda I)
&=(2-\lambda)\det\begin{pmatrix}-\lambda&-8\\4&12-\lambda\end{pmatrix}\\
&=(2-\lambda)(\lambda^2-12\lambda+32)\\
&=(2-\lambda)(\lambda-4)(\lambda-8).
\end{align*}
So the eigenvalues are
\[
\lambda_1=2,
\qquad
\lambda_2=4,
\qquad
\lambda_3=8.
\]

\subsubsection*{(ii) Eigenspaces and diagonalizability}
For $\lambda=2$,
\[
V_2^B=\ker(B-2I)=\ker\begin{pmatrix}
-2&0&-8\\
0&0&0\\
4&0&10
\end{pmatrix}.
\]
The equations give $x=0$ and $z=0$, so
\[
V_2^B=\Span\left\{\begin{pmatrix}0\\1\\0\end{pmatrix}\right\}.
\]
For $\lambda=4$,
\[
V_4^B=\ker(B-4I)=\ker\begin{pmatrix}
-4&0&-8\\
0&-2&0\\
4&0&8
\end{pmatrix}.
\]
Therefore
\[
x+2z=0,
\qquad y=0.
\]
Thus
\[
V_4^B=\Span\left\{\begin{pmatrix}-2\\0\\1\end{pmatrix}\right\}.
\]
For $\lambda=8$,
\[
V_8^B=\ker(B-8I)=\ker\begin{pmatrix}
-8&0&-8\\
0&-6&0\\
4&0&4
\end{pmatrix}.
\]
We get
\[
x+z=0,
\qquad y=0.
\]
Hence
\[
V_8^B=\Span\left\{\begin{pmatrix}1\\0\\-1\end{pmatrix}\right\}.
\]
All algebraic multiplicities and geometric multiplicities are equal to $1$, hence $B$ is diagonalizable.

\subsubsection*{(iii) Computing $B^n$}
Let
\[
D=\begin{pmatrix}
2&0&0\\
0&4&0\\
0&0&8
\end{pmatrix},
\qquad
P=\begin{pmatrix}
0&-2&1\\
1&0&0\\
0&1&-1
\end{pmatrix}.
\]
Then
\[
P^{-1}=\begin{pmatrix}
0&1&0\\
-1&0&-1\\
-1&0&-2
\end{pmatrix}.
\]
Since $B=PDP^{-1}$, we get
\[
B^n=PD^nP^{-1}.
\]
Therefore
\[
B^n=
\begin{pmatrix}
2\cdot4^n-8^n&0&2\cdot4^n-2\cdot8^n\\
0&2^n&0\\
8^n-4^n&0&2\cdot8^n-4^n
\end{pmatrix}.
\]

\newpage
\section*{Exercise 5.2}

\subsection*{Statement summary}
Over $\K=\C$, let $A\in\C^{n\times n}$ be the upper triangular matrix whose diagonal entries are all $i$ and whose first superdiagonal entries are also $i$, where $i^2=-1$. Determine the eigenvalues and eigenspaces of $A$, and decide for which $n$ the matrix is diagonalizable.

\subsection*{Solution}
Since the matrix is upper triangular,
\[
\det(A-\lambda I)=(i-\lambda)^n.
\]
Thus
\[
\lambda=i
\]
is the only eigenvalue.

Now
\[
V_i^A=\ker(A-iI).
\]
But
\[
A-iI=\begin{pmatrix}
0&i&0&\cdots&0&0\\
0&0&i&\cdots&0&0\\
0&0&0&\ddots&0&0\\
\vdots&\vdots&\vdots&\ddots&i&0\\
0&0&0&\cdots&0&i\\
0&0&0&\cdots&0&0
\end{pmatrix}.
\]
Therefore
\[
\begin{pmatrix}
0&i&0&\cdots&0&0\\
0&0&i&\cdots&0&0\\
0&0&0&\ddots&0&0\\
\vdots&\vdots&\vdots&\ddots&i&0\\
0&0&0&\cdots&0&i\\
0&0&0&\cdots&0&0
\end{pmatrix}
\begin{pmatrix}
x_1\\x_2\\\vdots\\x_n
\end{pmatrix}=0.
\]
This gives
\[
ix_j=0\quad\Rightarrow\quad x_j=0,
\qquad j=2,\dots,n.
\]
Hence
\[
V_i^A=\Span\left\{\begin{pmatrix}1\\0\\\vdots\\0\end{pmatrix}\right\}.
\]
So
\[
\dim V_i^A=1,
\qquad
m_i=n.
\]
Therefore $A$ is diagonalizable if and only if $n=1$.

\end{document}
